\subsection{拓扑群的完备化}

一致空间 \(G_{d}\) 可以视为它的完备化 \(\hat{G}_{d}\) 的一个稠密子空间。我们将研究 G 能否被视为一个完备 Hausdorff 群 \(G'\) 的一个稠密子群。若能，则一致空间 \(G_{d}'\) 必同构于 \(\hat{G}_{d}\)（第二章 § 3 第 6 目定理 2 的推论），因而我们必须能在 \(\hat{G}_{d}\) 上定义一个拓扑群结构，它在 G 上诱导给定的拓扑群结构。因此我们必须考察：1) 能否把函数 xy 与 \(x^{-1}\) 分别连续延拓到 \(\hat{G}_{d} \times \hat{G}_{d}\) 与 \(\hat{G}_{d}\) 上；2) 这样延拓后的函数是否确实在 \(\hat{G}_{d}\) 上定义一个群结构（这时它们必然在 \(\hat{G}_{d}\) 上定义一个在 G 上诱导给定结构的拓扑群结构）。接着我们要确立：3) 当前面的运算可行时，它们所定义的拓扑群是完备的。最后，我们将看到：4) 若存在一个满足给定条件的完备群，则它在同构意义下是唯一的。

\begin{enumerate}
\def\labelenumi{\arabic{enumi})}
\tightlist
\item
  把 \(xy\) 与 \(x^{-1}\) 连续延拓。由于函数 \(xy\) 与 \(x^{-1}\) 一般不是一致连续的，我们不能应用一致连续函数的延拓定理（第二章 § 3 第 6 目定理 2）。尽管如此，凭第二章 § 3 第 6 目命题 11 以及下面的命题，我们可以延拓 \(xy\)：
\end{enumerate}

命题 6 设 \(\mathfrak{F}\) 与 \(\mathfrak{G}\) 是 \(\mathbf{G}_d\) 上的两个 Cauchy 滤子。则滤子 \(\mathfrak{F} \times \mathfrak{G}\) 在映射 \((x, y) \to xy\) 下的像是 \(\mathbf{G}_d\) 上的一个 Cauchy 滤子基。让我们估计 \(xy\) 与 \(x'y'\) 在 \(\mathbf{G}_d\) 中的“接近程度”；换言之，让我们作乘积 \((x'y')(xy)^{-1} = x'y'y^{-1}x^{-1}\)。对每个 \(a \in \mathbf{G}\)，我们还可以写 \((x'y')(xy)^{-1} = (x'a^{-1})(ay'y^{-1}a^{-1})(ax^{-1})\)。我们将看到，通过适当选取 \(a\)，只要偶 \((x, y)\) 与 \((x', y')\) 属于 \(\mathfrak{F} \times \mathfrak{G}\) 的一个充分小的集，这个乘积的三个因子中的每一个都很小。设 V 是 \(e\) 在 G 中的任一邻域；则存在一个 \(\mathbf{V}_d\)-小集 A ∈ \(\mathfrak{F}\)。在 A 中选取 \(a\)，则若 \(x\) 与 \(x'\) 是 A 的任意两点，我们有 \(x'a^{-1} \in V\) 与 \(ax^{-1} \in V\)。另一方面，关系 \(ay'y^{-1}a^{-1} \in V\) 等价于

\[
y ^ {\prime} y ^ {- 1} \in a ^ {- 1} \mathrm{V}a = \mathrm{W},
\]

而由于 W 是 e 的一个邻域，存在一个 \(W_{d}\)-小集 \(B \in \mathfrak{G}\)。于是对 A × B 中所有 \((x, y)\) 与 \((x', y')\)，我们有 \((x'y')(xy)^{-1} \in V^{3}\)，证毕。

为使 \(x^{-1}\) 可以连续延拓到 \(\hat{G}_d\) 上，充要条件是：在对称映射 \(x \to x^{-1}\) 下，一个 Cauchy 滤子（\(G_d\) 上的）的像是 \(G_d\) 上的一个 Cauchy 滤子（第二章 § 3 第 6 目命题 11）。存在一些拓扑群，这个条件在其中不满足（参看第十章 § 3 习题 16）；在本证明的余下部分我们将假设它满足。

\begin{enumerate}
\def\labelenumi{\arabic{enumi})}
\setcounter{enumi}{1}
\item
  延拓后的函数 \(xy\) 与 \(x^{-1}\) 在 \(\hat{G}_d\) 上定义一个群结构。因为若把恒等延拓原理（第一章 § 8 第 1 目命题 2 推论 1）应用于函数 \(x(yz)\) 与 \((xy)z\)（它们定义在 \(\hat{G}_d \times \hat{G}_d \times \hat{G}_d\) 上且在稠密子空间 \(G_d \times G_d \times G_d\) 上相等），我们就看出合成法则 \((x, y) \to xy\) 在 \(\hat{G}_d\) 上是结合的。出于同样的理由，函数 \(x, ex, xe\) 在 \(\hat{G}_d\) 上相同，且函数 \(e, xx^{-1}, x^{-1}x\) 在 \(\hat{G}_d\) 上相同。
\item
  拓扑群 \(\hat{\mathbf{G}}_d\) 是完备的。设 \(\mathfrak{U}_d\) 是它的右一致结构，\(\mathfrak{U}\) 是 \(\hat{\mathbf{G}}_d\) 上的由 \(\mathbf{G}\) 的右一致结构的完备化所得的一致结构。则 \(\mathfrak{U}\) 与 \(\mathfrak{U}_d\) 在 \(\mathbf{G}\) 上诱导相同的一致结构，因而每个 Cauchy 滤子基 \(\mathfrak{B}\)（\(\mathbf{G}\) 上的，关于 \(\mathfrak{U}_d\) 的）也是关于 \(\mathfrak{U}\) 的一个 Cauchy 滤子基。现在 \(\mathfrak{B}\) 在 \(\hat{\mathbf{G}}_d\) 中收敛，因为 \(\mathcal{U}\) 是一个完备一致结构；且由于 \(\mathcal{U}\) 与 \(\mathcal{U}_d\) 在 \(\hat{\mathbf{G}}_d\) 上诱导相同的拓扑，可知（第二章 § 3 第 4 目命题 9）\(\mathcal{U}_d\) 是一个完备一致结构。这个结论还表明 \(\mathcal{U}\) 与 \(\mathcal{U}_d\) 相同（第二章 § 3 第 6 目定理 2 的推论）。
\item
  唯一性。这由第 3 目命题 5 得出。
\end{enumerate}

概括起来，我们证明了下面的定理：

定理 1 Hausdorff 拓扑群 G 同构于一个完备群 \(\hat{G}\) 的一个稠密子群，当且仅当关于 G 的右一致结构的一个 Cauchy 滤子在对称映射 \(x \to x^{-1}\) 下的像是关于这个一致结构的一个 Cauchy 滤子。完备群 \(\hat{G}\)（它称为 G 的完备化）这时是唯一的（在同构意义下）。

命题 7 设 G 是一个具有完备化 \(\hat{G}\) 的 Hausdorff 拓扑群。则 G 中单位元的诸邻域在 \(\hat{G}\) 中的闭包构成 \(\hat{G}\) 中单位元的一个基本邻域系。

由于 \(\hat{G}\) 是正则的，\(\hat{G}\) 中单位元的每个邻域含有 e 在 \(\hat{G}\) 中的一个开邻域 U 的闭包 V，而 V 也是 U 在 G 上的迹的闭包。

设 G 是一个未必 Hausdorff 的群；设 \(N = \overline{e}\)，\(G' = G/N\) 是伴随于 G 的 Hausdorff 群（§ 2 第 6 目）。若 \(G'\) 有一个完备化 \(\hat{G}'\)，这个完备化称为 G 的 Hausdorff 完备化，记作 \(\hat{G}\)；这时 \(\hat{G}_{d}'\)（相应地 \(\hat{G}_{s}'\)）是一致空间 \(G_{d}\)（相应地 \(G_{s}\)）的 Hausdorff 完备化（第二章 § 3 第 7 目）。

命题 8 设 G 是一个具有 Hausdorff 完备化 \(\hat{G}'\) 的拓扑群。则 G 到一个完备 Hausdorff 群 H 中的每个连续同态 \(u\) 可以唯一地分解为 \(u = v \circ \varphi\)，其中 \(v\) 是 \(\hat{G}'\) 到 H 中的一个连续同态，而 \(\varphi\) 是 G 到 \(\hat{G}'\) 中的典范映射（G' 到 \(\hat{G}'\) 中的典范单射与 G 到 G/N = G' 上的典范同态 \(\psi\) 的复合）。

由于 \(u\) 的核是闭的且含有 \(e\)，它含有 \(N\)，因而 \(u\) 可以写作 \(u = w \circ \psi\)，其中 \(w\) 是 \(G'\) 到 \(H\) 中的一个连续同态；现在对 \(w\) 应用第 3 目命题 5。

